% Source: https://artofproblemsolving.com/community/c3h1675627_let_pn_denote_the_number_of_partitions_of_n_and_let_pn_m_denote_the_number
% Stable identity: aops-topic-1675627-q01
Write a partition of $2n$ into exactly $n$ parts as $\lambda_1\ge\lambda_2\ge\cdots\ge\lambda_n\ge1$ with $\lambda_1+\cdots+\lambda_n=2n$. Subtract $1$ from every part:
\[
\mu_i=\lambda_i-1\ge0,\qquad \mu_1\ge\mu_2\ge\cdots\ge\mu_n\ge0,\qquad \mu_1+\cdots+\mu_n=2n-n=n .
\]
Deleting the zero parts, $(\mu_i)$ is a partition of $n$ with at most $n$ parts. Conversely, any partition of $n$ has at most $n$ parts (each part is at least $1$); padding it with zeros to exactly $n$ entries and adding $1$ to each entry gives a partition of $2n$ into exactly $n$ parts. The two maps are inverse to each other, so they form a bijection between the partitions of $2n$ into exactly $n$ parts and all partitions of $n$. Hence $p(2n,n)=p(n)$. $\blacksquare$

For example, $n=3$: the partitions $3$, $2+1$, $1+1+1$ of $3$ correspond to $4+1+1$, $3+2+1$, $2+2+2$, the partitions of $6$ into exactly $3$ parts.
